%============================================================
% CHAPTER 3 -- SOFTWARE DESIGN DESCRIPTION (SDD)
%============================================================

\chapter{Software Design Description (SDD)}

%------------------------------------------------------------
\section{Introduction}
%------------------------------------------------------------

\subsection{Design Overview}

This Software Design Description (SDD) translates the functional and non-functional
requirements from the SRS into concrete structural decisions: how the system is split into
layers and components, how those components interact at runtime, and how persistent data
is organised. The design choices here are not arbitrary --- each one maps back to one or
more requirements, and the traceability matrix in Section~3.1.2 makes those mappings
explicit. The chapter closes with a class diagram covering the key domain entities and
the backend service layer.

MediLingo AI uses a \textbf{three-tier client--server architecture}: a React single-page
application running in the browser, a stateless FastAPI application server, and a dual
persistence layer made up of MongoDB for user and session data and ChromaDB for the indexed
medical knowledge base. The three tiers communicate over HTTPS with JSON payloads; refresh
tokens travel as HTTP-only cookies, keeping them away from JavaScript running in the page.
No business logic lives in the frontend --- the SPA is purely responsible for rendering
state it receives from the API.

\subsection{Requirements Traceability Matrix}

Table~\ref{tab:rtm} maps each functional requirement from the SRS to the design component
responsible for satisfying it, along with the section of this chapter where that component
is described in detail. Any requirement missing a matching row would be a gap in the design.

\begin{longtable}{|>{\bfseries}p{1.8cm}|p{5.2cm}|p{3.8cm}|p{2.0cm}|}
  \caption{Requirements Traceability Matrix}\label{tab:rtm}\\
  \hline
  \theader{Req.\ ID} & \theader{Requirement Summary} & \theader{Design Component} & \theader{Section} \\
  \hline
  \endfirsthead
  \multicolumn{4}{l}{\small\itshape\color{darkgray}(Table \ref{tab:rtm} continued)}\\
  \hline
  \theader{Req.\ ID} & \theader{Requirement Summary} & \theader{Design Component} & \theader{Section} \\
  \hline
  \endhead
  \hline
  \endfoot
  FR-01 & User Registration & Auth Router, UserDoc Model & \ref{sec:auth-router} \\
  \hline
  FR-02 & User Login & Auth Router, JWT Service & \ref{sec:auth-router} \\
  \hline
  FR-03 & Session Persistence (JWT) & Auth Router, apiClient & \ref{sec:auth-router} \\
  \hline
  FR-04 & Logout & Auth Router & \ref{sec:auth-router} \\
  \hline
  FR-05 & Submit Medical Query & Query Router, RAG Engine & \ref{sec:query-router} \\
  \hline
  FR-06 & Language Selection & Frontend Language Store, RAG Engine Prefix Stripping & \ref{sec:frontend} \\
  \hline
  FR-07 & Multilingual Response & RAG Engine, LLM System Prompt & \ref{sec:rag-engine} \\
  \hline
  FR-08 & Retrieve Context Chunks & RAG Engine, Vector Store & \ref{sec:rag-engine} \\
  \hline
  FR-09 & Create Chat Session & Sessions Router, SessionDoc & \ref{sec:sessions-router} \\
  \hline
  FR-10 & Persist Chat History & Sessions Router, MessageDoc & \ref{sec:sessions-router} \\
  \hline
  FR-11 & Load Previous Sessions & Sessions Router & \ref{sec:sessions-router} \\
  \hline
  FR-12 & Delete Session & Sessions Router & \ref{sec:sessions-router} \\
  \hline
  FR-13 & Voice Input & useVoiceInput Hook (Web Speech API) & \ref{sec:frontend} \\
  \hline
  FR-14 & Voice Output & useVoiceOutput Hook (Web Speech API) & \ref{sec:frontend} \\
  \hline
  FR-15 & Admin Login & Auth Router (role check) & \ref{sec:auth-router} \\
  \hline
  FR-16 & PDF Ingestion (Admin) & Pipeline Router, RAG Ingestion Pipeline & \ref{sec:pipeline-router} \\
  \hline
\end{longtable}

%------------------------------------------------------------
\section{System Architectural Design}
%------------------------------------------------------------

\subsection{Chosen System Architecture}

MediLingo AI uses a \textbf{four-tier layered architecture}:

\begin{itemize}
  \item \textbf{Presentation Tier} --- React 19 SPA compiled with Vite, served as static
        files. Talks to the application tier through REST only; no business logic lives here.
        UI state is managed by Zustand stores and voice I/O goes through the browser's
        Web Speech API.

  \item \textbf{Application Tier} --- FastAPI server. Validates requests, enforces
        authentication, coordinates calls to the AI/ML and persistence tiers, and returns
        structured JSON. Stateless by design, which means scaling horizontally is
        straightforward in principle.

  \item \textbf{AI / ML Processing Tier} --- The RAG engine, initialised once at server
        startup and held as a singleton. Handles embedding generation
        (sentence-transformers), vector similarity search (ChromaDB), language-prefix
        stripping, prompt assembly, and LLM invocation (Groq API, LLaMA 3.3-70b). This
        is the most computationally intensive part of the stack.

  \item \textbf{Persistence Tier} --- Two stores: MongoDB (Motor async driver) for user
        accounts, sessions and messages; ChromaDB (local persistence) for the encoded
        medical knowledge base. Both are initialised once at startup via the FastAPI
        lifespan handler.
\end{itemize}

Figure~\ref{fig:arch} presents the system architectural diagram.

% ── System Architectural Diagram ──────────────────────────
\begin{figure}[H]
  \centering
  \includegraphics[width=\textwidth]{architecture}
  \caption{MediLingo AI --- Three-Tier System Architecture}
  \label{fig:arch}
\end{figure}

\subsection{Discussion of Alternative Designs}

\subsubsection{Monolithic vs.\ Microservices}

A microservices architecture was evaluated --- separate deployable units for auth, chat,
ingestion and retrieval. It was rejected. The overhead of inter-service communication,
distributed tracing, and individual deployment pipelines is real, and the benefit at this
scale is not. A single developer building an academic prototype does not need Kubernetes.
The layered monolith achieves the same separation of concerns through module boundaries
rather than process boundaries, which is much easier to develop, debug and test.

\subsubsection{Server-Side Rendering vs.\ SPA}

Next.js with server-side rendering was considered as an alternative to the Vite SPA. SSR
improves initial page load and SEO, which matters for public-facing content. MediLingo AI's
chat interface is inherently interactive and session-gated --- there is no content for
search engines to index, no cold visitors landing on a page. The Vite build was preferred:
faster development iteration, simpler deployment as static files, and cleaner integration
with Zustand for keeping access tokens in memory rather than exposing them through
server-rendered HTML.

\subsubsection{Relational vs.\ Document Store}

PostgreSQL with pgvector was a tempting option --- one database handling both relational
user data and vector search. ChromaDB plus MongoDB was chosen instead. ChromaDB's
purpose-built HNSW index and embedding API make the retrieval path significantly
simpler, and MongoDB's flexible schema means the message and session structures can
evolve without writing and running migrations every time something changes. For an
academic project where the schema is still being iterated on, that flexibility mattered.

\subsection{System Interface Description}

\begin{itemize}
  \item \textbf{User Interface} — Browser-based React SPA. Input: text box, microphone
        button, language selector. Output: streamed chat bubbles, text-to-speech playback.

  \item \textbf{API Interface} — REST over HTTP/HTTPS. Base URL
        \texttt{/api/v1}. Authentication via \texttt{Authorization: Bearer <access\_token>}
        header. Refresh via HTTP-only \texttt{refresh\_token} cookie sent automatically
        by the browser.

  \item \textbf{LLM Interface} — Groq REST API. The RAG engine sends a
        \texttt{chat.completions} request containing a system prompt (with language
        instruction and retrieved context) and the user message. Model:
        \texttt{llama-3.3-70b-versatile}.

  \item \textbf{Database Interface} — Motor (async PyMongo) for MongoDB;
        ChromaDB Python client for vector operations. Both connections are initialised
        once at server startup via the FastAPI lifespan handler.
\end{itemize}

%------------------------------------------------------------
\section{Detailed Description of Components}
\label{sec:components}
%------------------------------------------------------------

\subsection{Frontend (React SPA)}
\label{sec:frontend}

The frontend is a React 19 / TypeScript application built with Vite. It is structured
around pages, Zustand stores, and custom hooks that keep side-effectful code away from
the rendering layer. Most components are stateless; they display data received through
props or from the stores, and they call hook functions to trigger any side effects.

\begin{itemize}
  \item \textbf{router.tsx} --- All client-side routes, defined with React Router v6.
        Protected routes check the auth store and redirect to \texttt{/login} when no
        valid access token is present.

  \item \textbf{App.tsx} --- The root component. On mount it silently calls
        \texttt{/auth/refresh}; if the refresh cookie is still valid the access token
        is restored and the user lands directly on the chat page without re-entering
        credentials. This is how the ``stay logged in'' experience works.

  \item \textbf{useAuthStore (Zustand)} --- Holds \texttt{user} and
        \texttt{accessToken} in memory only. Tokens are never written to
        \texttt{localStorage} or \texttt{sessionStorage} --- they disappear when the
        tab closes.

  \item \textbf{apiClient.ts} --- All API calls go through here. It wraps \texttt{fetch}
        with a configurable timeout and handles 401 responses automatically: calls
        \texttt{/auth/refresh}, updates the store, and retries the original request once.
        If the retry also fails, it clears auth state and redirects to \texttt{/login}.

  \item \textbf{useVoiceInput / useVoiceOutput} --- Thin wrappers around
        \texttt{SpeechRecognition} and \texttt{SpeechSynthesis}. The output hook picks
        a locale-appropriate voice based on the active language setting.
\end{itemize}

\subsection{Authentication Router}
\label{sec:auth-router}

\texttt{routers/auth.py} exposes five endpoints under \texttt{/auth}:

\begin{itemize}
  \item \textbf{POST /register} --- Validates uniqueness of both username and email, hashes
        the password with bcrypt via \texttt{passlib}, stores a \texttt{UserDoc} in MongoDB,
        and returns an access/refresh token pair. A duplicate email gets HTTP~400.

  \item \textbf{POST /login} --- Verifies credentials against the stored bcrypt hash and
        issues a new token pair. The refresh token goes into an HTTP-only, Secure,
        SameSite=Lax cookie with a 7-day TTL.

  \item \textbf{POST /refresh} --- Reads the refresh cookie, verifies the JWT signature and
        expiry, then issues a new 15-minute access token. The remaining TTL on the refresh
        token is preserved, not reset.

  \item \textbf{POST /logout} --- Clears the refresh cookie. Nothing else needs to happen
        server-side because access tokens are not persisted anywhere.

  \item \textbf{GET /me} --- Returns the current user's profile using the access token from
        the Authorization header.
\end{itemize}

All tokens are signed HS256 JWTs. The \texttt{SECRET\_KEY} used for signing is loaded
from the environment at startup and is never present in the source code.

\subsection{Sessions Router}
\label{sec:sessions-router}

\texttt{routers/sessions.py} manages chat sessions and their messages under
\texttt{/sessions}. Every endpoint here requires a valid access token.

\begin{itemize}
  \item \textbf{POST /sessions} --- Creates a new \texttt{SessionDoc} in MongoDB with a
        UUID, the authenticated user's ID, a display title and an empty message list.

  \item \textbf{GET /sessions} --- Returns the user's sessions, sorted by last-updated
        descending. This is what populates the sidebar.

  \item \textbf{GET /sessions/\{id\}} --- Returns a single session with its full message
        history. Called when a user clicks on a session to resume it.

  \item \textbf{POST /sessions/\{id\}/messages} --- Appends a \texttt{MessageDoc} (role,
        content, timestamp, language) to the session. Called by the query router after
        each exchange, not directly by the user.

  \item \textbf{DELETE /sessions/\{id\}} --- Deletes the session and all its embedded
        messages.
\end{itemize}

\subsection{Query and Pipeline Routers}
\label{sec:query-router}
\label{sec:pipeline-router}

\texttt{routers/query.py} handles all AI-related endpoints:

\begin{itemize}
  \item \textbf{POST /query} — The primary chat endpoint. Accepts \texttt{query} (the
        full language-prefixed string), strips the prefix, embeds the clean query,
        retrieves the top-\textit{k} context chunks from ChromaDB, assembles the Groq
        prompt, and streams the LLM response back to the caller.

  \item \textbf{POST /retrieve} — Returns only the raw retrieved context chunks, without
        LLM generation. Used for debugging and evaluation.

  \item \textbf{GET /config} — Returns current RAG configuration (chunk size, overlap,
        top-k, model name).

  \item \textbf{GET /info/\{collection\}} — Returns metadata about a ChromaDB collection
        (document count, embedding dimension).
\end{itemize}

\texttt{routers/query.py} also contains pipeline management endpoints under
\texttt{/pipeline}:

\begin{itemize}
  \item \textbf{POST /pipeline/ingest-sync} — Synchronously ingests all PDF files
        found in \texttt{backend/datasets/}. Admin-only. Returns a summary of documents
        indexed, chunks created, and time elapsed.

  \item \textbf{POST /pipeline/ingest} — Same as above but runs as a FastAPI
        \texttt{BackgroundTask}, returning immediately with a task-accepted response.
\end{itemize}

\subsection{RAG Engine}
\label{sec:rag-engine}

The RAG engine (\texttt{rag/engine.py}) is the core AI component of the system.
It is initialised once at server startup through \texttt{services/rag\_service.py} and
held as a module-level singleton. The processing pipeline for each incoming query has
five steps:

\begin{enumerate}
  \item \textbf{Prefix Stripping} --- The language instruction prefix (e.g.,
        \textit{``Respond in Urdu: ''}) is detected and stripped before anything else
        happens. This step is critical for retrieval quality: the clean medical question
        is what gets embedded and searched, not the prefixed string.

  \item \textbf{Embedding} --- The clean query is passed to
        \texttt{EmbeddingGenerator.encode()}, which runs \texttt{all-MiniLM-L6-v2} via
        \texttt{sentence-transformers} to produce a 384-dimensional dense vector.

  \item \textbf{Vector Retrieval} --- The 384-dim query vector is submitted to
        \texttt{ChromaDBVectorStore.query()}, which returns the top-$k$ document chunks
        by cosine similarity.

  \item \textbf{Prompt Assembly} --- A structured Groq chat-completions prompt is built:
        the system message contains the language instruction, a medical-domain persona
        and the retrieved context passages; the user message contains the original
        prefixed query.

  \item \textbf{LLM Invocation} --- The assembled prompt is sent to the Groq API
        (\texttt{llama-3.3-70b-versatile}). The response text is returned to the calling
        router, which saves it to MongoDB and returns it to the frontend.
\end{enumerate}

\subsection{Ingestion Pipeline}
\label{sec:ingestion}

\texttt{rag/pipeline.py} contains \texttt{RAGIngestionPipeline}, which handles the
one-time (or on-demand) process of building the medical knowledge base:

\begin{enumerate}
  \item \textbf{PDF Loading} --- \texttt{rag/loader.py} uses \texttt{pdfplumber} as the
        primary extractor, falling back to \texttt{pypdf} for files that pdfplumber
        cannot handle cleanly. Every PDF in \texttt{backend/datasets/} is processed.

  \item \textbf{Text Cleaning} --- \texttt{rag/preprocessor.py} normalises whitespace,
        expands common medical abbreviations, strips headers and footers, and removes
        non-printable characters. The goal is clean, consistent text going into the
        embedder.

  \item \textbf{Semantic Chunking} --- \texttt{rag/chunker.py} splits cleaned text into
        fixed-size chunks of 1000 characters with 200-character overlap. The overlap
        ensures sentences split across a chunk boundary still appear complete in at
        least one chunk. Short chunks below a minimum length are discarded.

  \item \textbf{Embedding and Indexing} --- Each chunk is embedded and upserted into
        the ChromaDB collection with source filename and chunk index as metadata.
        Upsert rather than insert means re-running the pipeline on an already-indexed
        document does not create duplicates.
\end{enumerate}

\subsection{Data Models}
\label{sec:models}

Three MongoDB document classes are defined in \texttt{models/}:

\begin{itemize}
  \item \textbf{UserDoc} --- \texttt{id} (UUID4 str), \texttt{username}, \texttt{email},
        \texttt{hashed\_password}, \texttt{role} (\textit{user} or \textit{admin}),
        \texttt{created\_at}. Unique indexes on \texttt{username} and \texttt{email} are
        created at startup to enforce deduplication at the database level.

  \item \textbf{SessionDoc} --- \texttt{id} (UUID4 str), \texttt{user\_id},
        \texttt{title}, \texttt{created\_at}, \texttt{updated\_at}, and
        \texttt{messages} --- an embedded array of \texttt{MessageDoc}. The index on
        \texttt{user\_id} makes the ``fetch my sessions'' query fast regardless of
        collection size.

  \item \textbf{MessageDoc} --- \texttt{id} (UUID4 str), \texttt{role} (\textit{user}
        or \textit{assistant}), \texttt{content}, \texttt{language}, \texttt{timestamp}.
        Stored embedded inside \texttt{SessionDoc} rather than in a separate collection.
\end{itemize}

All three models expose \texttt{to\_mongo()} and \texttt{from\_mongo()} helpers.
UUID4 string IDs are used throughout to decouple application identity from MongoDB's
internal \texttt{\_id} field --- application code never touches \texttt{\_id} directly.

%------------------------------------------------------------
\section{Architectural Diagram}
%------------------------------------------------------------

Figure~\ref{fig:arch} in Section~3.2.1 shows the four-tier layered architecture.
Figure~\ref{fig:rag-flow} below zooms in on the AI/ML tier and traces the full
query processing flow from the moment a prefixed query arrives at the RAG engine to
the moment a response is returned to the router.

\begin{figure}[H]
  \centering
  \includegraphics[width=0.55\textwidth]{rag_flow}
  \caption{RAG Query Processing Flow}
  \label{fig:rag-flow}
\end{figure}

%------------------------------------------------------------
\section{Sequence Diagrams}
\label{sec:sequence-diagrams}
%------------------------------------------------------------

This section presents sequence diagrams for each of the ten use cases identified in the
SRS. All ten diagrams are included as exported PNG sequence diagrams.

\subsection{SSD-1: Register Account}

\begin{figure}[H]
  \centering
  \includegraphics[width=0.85\textwidth]{ssd1_register}
  \caption{SSD-1: Register Account}
  \label{fig:ch3-ssd1}
\end{figure}

\subsection{SSD-2: Login}

\begin{figure}[H]
  \centering
  \includegraphics[width=0.85\textwidth]{ssd2_login}
  \caption{SSD-2: Login}
  \label{fig:ch3-ssd2}
\end{figure}

\subsection{SSD-3: Submit Medical Query}

\begin{figure}[H]
  \centering
  \includegraphics[width=\textwidth]{ssd3_submit_query}
  \caption{SSD-3: Submit Medical Query}
  \label{fig:ch3-ssd3}
\end{figure}

\subsection{SSD-4: Voice Input}

\begin{figure}[H]
  \centering
  \includegraphics[width=0.85\textwidth]{ssd5_voice_input}
  \caption{SSD-4: Voice Input --- Speech Recognition Flow}
  \label{fig:ch3-ssd4}
\end{figure}

\subsection{SSD-5: Voice Output}

\begin{figure}[H]
  \centering
  \includegraphics[width=0.85\textwidth]{ssd6_voice_output}
  \caption{SSD-5: Voice Output --- Text-to-Speech Playback Flow}
  \label{fig:ch3-ssd5}
\end{figure}

\subsection{SSD-6: Create New Session}

\begin{figure}[H]
  \centering
  \includegraphics[width=0.85\textwidth]{ssd4_sessions}
  \caption{SSD-6: Create New Session}
  \label{fig:ch3-ssd6}
\end{figure}

\subsection{SSD-7: Load Previous Sessions}

\begin{figure}[H]
  \centering
  \includegraphics[width=0.90\textwidth]{ssd4_sessions}
  \caption{SSD-7 \& SSD-8: Session Management --- Create, Load, and Delete}
  \label{fig:ch3-ssd7}
\end{figure}

\subsection{SSD-8: Delete Session}

The delete session flow is covered in Figure~\ref{fig:ch3-ssd7} above, which shows all
three session management operations (create, load, and delete) in a single combined
sequence diagram.

\subsection{SSD-9: Admin Login}

\begin{figure}[H]
  \centering
  \includegraphics[width=0.85\textwidth]{ssd2_login}
  \caption{SSD-9: Admin Login --- same auth flow as SSD-2, role claim verified server-side}
  \label{fig:ch3-ssd9}
\end{figure}

\subsection{SSD-10: Ingest PDF Documents}

\begin{figure}[H]
  \centering
  \includegraphics[width=0.90\textwidth]{ssd10_ingest}
  \caption{SSD-10: Ingest PDF Documents}
  \label{fig:ch3-ssd10}
\end{figure}

%------------------------------------------------------------
\section{Class Diagram}
%------------------------------------------------------------

Figure~\ref{fig:class} shows the key domain classes, their attributes and their
relationships. Frontend Zustand stores are left out --- this diagram focuses on the
backend data model and the service layer that drives it. The full picture would include
the React hooks and stores, but adding them here would make the diagram unreadable
without adding much that is not already covered in the frontend component description.

\begin{figure}[H]
  \centering
  \includegraphics[width=\textwidth]{class_diagram}
  \caption{MediLingo AI --- Class Diagram}
  \label{fig:class}
\end{figure}
